\documentclass[11pt]{article}
\usepackage[margin=1in]{geometry}
\usepackage{booktabs}
\usepackage{amsmath}
\usepackage{hyperref}
\usepackage{authblk}

\title{From Pinning to Amplification:\\Evidence of a Regime Shift in S\&P~500 Options Expiration Dynamics,\\2016--2025}

\author{Nathan Elms}
\affil{Independent Researcher}

\date{April 2026}

\begin{document}

\maketitle

\begin{abstract}
Prior work has established that options open interest can either pin or amplify underlying price movement near expiration, depending on the sign of aggregate dealer gamma exposure (Avellaneda \& Lipkin 2003; Jeannin et~al.\ 2008; Barbon \& Buraschi 2021). Golez \& Jackwerth (2012) documented pinning in S\&P~500 futures through 2009. We extend their analysis to 2016--2025, a period marked by the proliferation of weekly and zero-days-to-expiration (0DTE) options, using 1.6~million near-expiry SPY options records matched with E-mini S\&P~500 futures at one-minute resolution across 2{,}294 trading days. Five tests find no evidence of pinning in this modern sample. The single statistically significant result ($p < 0.001$) shows that days with the highest ATM open interest exhibit 16\% wider daily ranges than days with the lowest---consistent with the gamma amplification mechanism described by Barbon \& Buraschi (2021) rather than the pinning documented in earlier periods. These results suggest that S\&P~500 expiration dynamics may have shifted from a pinning regime to an amplification regime, potentially driven by structural changes in dealer positioning associated with short-dated options growth. This investigation was prompted by a structural profiling tool (Database Whisper) that identified open interest as more structurally informative than implied volatility in options data.
\end{abstract}

\section{Introduction}

Options pinning---the tendency of underlying prices to gravitate toward strikes with concentrated open interest near expiration---is among the best-documented market microstructure effects in equity options. The theoretical foundation was laid by Avellaneda \& Lipkin~\cite{avellaneda2003}, who derived the delta-hedging-driven drift that pulls prices toward high open interest strikes. Ni, Pearson, \& Poteshman~\cite{ni2005} provided definitive empirical evidence of pinning in U.S.\ equity options, estimating aggregate market capitalization effects of \$9~billion per expiration date. This work has been extended and refined by multiple authors~\cite{jeannin2008,filippou2022,chung2013}.

An important theoretical insight, however, is that the \textit{direction} of the hedging effect depends on dealer positioning. Jeannin et~al.~\cite{jeannin2008} showed that when market makers are \textit{long} options (positive gamma), their hedging activity dampens price movement---pinning. When they are \textit{short} options (negative gamma), their hedging \textit{amplifies} price movement. Barbon \& Buraschi~\cite{barbon2021} empirically documented this gamma fragility effect, showing that stocks with large negative dealer gamma exhibit intraday momentum rather than mean-reversion.

For the S\&P~500, the most comprehensive study is Golez \& Jackwerth~\cite{golez2012}, who found pinning in S\&P~500 futures on serial option expirations through November~2009. Their work predates several structural changes in the options market, including the explosive growth of weekly options (introduced broadly in 2010--2012) and zero-days-to-expiration (0DTE) contracts, which by 2023 accounted for approximately 50\% of daily SPX options volume.

These structural changes may have altered the balance of dealer positioning. If the proliferation of short-dated options has shifted dealers toward net short gamma positions more frequently, the theoretical framework of Jeannin et~al.\ and Barbon \& Buraschi predicts a shift from pinning to amplification.

We test this possibility using modern data (2016--2025) across five variations of the pinning hypothesis.

\subsection{Structural Motivation}

This investigation was prompted by an unexpected finding from Database Whisper~\cite{elms2026dw}, a structural profiling algorithm that automatically discovers which fields best distinguish records in structured datasets. Applied to 24.7~million SPY options records, the algorithm identified the following discriminator hierarchy among options sharing the same date, type, and expiration bucket:

\begin{center}
\texttt{delta} (72\%) $\to$ \texttt{open\_interest} (59\%) $\to$ \texttt{implied\_volatility} (52\%) $\to$ \texttt{volume} (43\%)
\end{center}

Open interest ranked as structurally more informative than implied volatility for distinguishing options---a finding that contradicts the typical practitioner focus on volatility over positioning. This structural prominence motivated a closer examination of open interest's empirical effects.

\section{Data}

\textbf{Options:} 24.7~million SPY options records, January~2008 to December~2025, containing strike, type, mark, bid, ask, volume, open interest, implied volatility, and full Greeks. Near-expiry subset ($\text{DTE} \leq 3$): 1{,}566{,}296 records.

\textbf{Underlying:} E-mini S\&P~500 futures (@ES) at one-minute resolution, October~2016 to April~2026 (3.3~million bars). Regular trading hours (9:30--16:00~ET): 937{,}182 bars across 2{,}436 trading days.

\textbf{Matching:} Each trading day's SPY options open interest profile is matched with the corresponding ES intraday price action (SPY $\approx$ ES/10 for strike-level comparisons). Merged dataset: 2{,}294 trading days.

\section{Five Tests}

For each trading day, we identify the SPY strike with maximum near-expiry ($\text{DTE} \leq 3$) open interest. We also compute total ATM open interest (options with $|\text{delta}| \in [0.35, 0.65]$). All tests examine whether price behavior relates to these open interest measures.

\subsection{Test 1: Close-vs-Open Distance}

Does price close nearer to the maximum-OI strike than it opened?

\textbf{Result:} Price moved toward the max-OI strike on 44.9\% of days and away on 55.1\%. Not significant ($p = 0.40$).

\subsection{Test 2: OI-Level Effect}

Does the gravitation effect strengthen with higher open interest?

\textbf{Result:} High-OI days moved toward the max-OI strike 47.5\% of the time vs.\ 42.2\% for low-OI days. Directionally consistent with pinning but not significant ($p = 0.19$).

\subsection{Test 3: Daily Range Compression}

Does concentrated open interest compress the daily trading range?

\begin{table}[h]
\centering
\begin{tabular}{lccc}
\toprule
OI Tercile & Mean Daily Range & Median & Days \\
\midrule
Low & 0.889\% & 0.688\% & 765 \\
Medium & 0.999\% & 0.791\% & 764 \\
High & 1.033\% & 0.847\% & 765 \\
\bottomrule
\end{tabular}
\caption{Daily range by near-expiry ATM open interest level. The relationship is monotonically positive.}
\label{tab:range}
\end{table}

\textbf{Result:} The opposite of pinning. High-OI days exhibit 16.2\% wider ranges than low-OI days. Statistically significant ($t = -3.67$, $p = 0.0003$). This is consistent with the gamma amplification mechanism of Barbon \& Buraschi~\cite{barbon2021}.

\subsection{Test 4: Intraday Time Near Max-OI Strike}

Does price spend more time near the max-OI strike on high-OI days?

\textbf{Result:} No meaningful difference detected on a 200-day sample.

\subsection{Test 5: Final-Hour Convergence}

Does price converge toward the max-OI strike in the final trading hour, when gamma effects should peak?

\textbf{Result:} Price converged from noon to close on 43.0\% of high-OI days vs.\ 45.0\% of low-OI days. No significant difference ($p = 0.78$).

\section{Summary of Results}

\begin{table}[h]
\centering
\begin{tabular}{llcl}
\toprule
Test & Metric & p-value & Finding \\
\midrule
V1: Close vs open & Toward: 44.9\% & 0.40 & No pinning \\
V2: OI-level effect & High OI toward: 47.5\% & 0.19 & No pinning \\
V3: Range compression & High OI: +16.2\% range & 0.0003 & Amplification \\
V4: Time near strike & No concentration & --- & No pinning \\
V5: Final-hour convergence & High OI: 43.0\% & 0.78 & No pinning \\
\bottomrule
\end{tabular}
\caption{Summary of five tests across 2{,}294 trading days.}
\label{tab:summary}
\end{table}

\section{Discussion}

\subsection{Reconciling with Prior Literature}

Our results do not contradict the pinning literature. Pinning has been convincingly demonstrated in individual equity options~\cite{ni2005,filippou2022}. Golez \& Jackwerth~\cite{golez2012} found pinning in S\&P~500 futures through 2009. Our finding is that, in the 2016--2025 sample, the S\&P~500 no longer exhibits pinning. The single significant result points in the opposite direction.

This is consistent with the theoretical framework. Jeannin et~al.~\cite{jeannin2008} showed that the direction of the hedging effect depends on the sign of dealer gamma. Barbon \& Buraschi~\cite{barbon2021} empirically confirmed that negative dealer gamma produces amplification. If the growth of short-dated options has shifted aggregate dealer positioning toward net short gamma more frequently, a transition from pinning to amplification is exactly what the theory predicts.

\subsection{Possible Structural Drivers}

Several market structure changes between the Golez \& Jackwerth sample (ending 2009) and ours (beginning 2016) may explain the shift:

\begin{itemize}
\item \textbf{0DTE proliferation:} By 2023, zero-days-to-expiration contracts represent approximately 50\% of daily SPX options volume. These contracts carry extreme gamma, and dealers who are short them experience amplified hedging demands.
\item \textbf{Weekly options growth:} Weekly expirations, introduced broadly in 2010--2012, distribute open interest across more frequent expirations, potentially changing the concentration dynamics that drive pinning.
\item \textbf{Retail options participation:} The growth of retail options trading since 2020 has altered the composition of open interest, with retail traders more likely to buy options (putting dealers short gamma) than to sell them.
\end{itemize}

We note these as potential explanations rather than confirmed mechanisms. Disentangling these effects would require dealer-level positioning data, which is not publicly available.

\subsection{Confounds}

An important confound is that high open interest near expiry is not exogenous. Traders accumulate positions in anticipation of scheduled events (FOMC, CPI, earnings), and those events independently cause large moves. The wider ranges on high-OI days may partly reflect this selection effect.

Additionally, our use of SPY options matched with ES futures introduces a cross-product approximation (SPY $\approx$ ES/10). While SPY and ES track closely, slight divergences may affect strike-level precision.

\subsection{Role of Structural Profiling}

The structural profiling tool that motivated this investigation (Database Whisper) identified open interest as more structurally informative than implied volatility in options data---a ranking that contradicts conventional practitioner wisdom. This structural finding did not predict the direction of the result. It identified that open interest carried more disambiguation power than expected, which prompted the empirical question. The data answered it.

This illustrates a potential workflow: automated structural analysis surfaces unexpected feature importance, which generates testable hypotheses, which produce empirical findings. The tool identified the question; the five tests answered it.

\section{Conclusion}

We find no evidence of options pinning in S\&P~500 expiration dynamics across five tests on 2{,}294 trading days (2016--2025). The significant finding ($p < 0.001$) shows the opposite: high near-expiry ATM open interest is associated with wider, not narrower, daily ranges. This is consistent with the gamma amplification mechanism documented by Barbon \& Buraschi~(2021) and the theoretical predictions of Jeannin et~al.~(2008) under net short dealer gamma.

We interpret these results as evidence of a possible regime shift in S\&P~500 expiration dynamics---from pinning in earlier periods (as documented by Golez \& Jackwerth through 2009) to amplification in the modern era, potentially driven by the structural growth of short-dated options and changes in dealer positioning. We offer this interpretation as a hypothesis for further investigation, not as a definitive conclusion.

These findings complement rather than contradict the existing pinning literature. Pinning remains well-documented in individual equities. The S\&P~500, however, may have outgrown it.

\begin{thebibliography}{99}

\bibitem{avellaneda2003}
M.~Avellaneda and M.~D.~Lipkin, ``A Market-Induced Mechanism for Stock Pinning,'' \textit{Quantitative Finance}, 3(6), 2003.

\bibitem{ni2005}
S.~X.~Ni, N.~D.~Pearson, A.~M.~Poteshman, ``Stock Price Clustering on Option Expiration Dates,'' \textit{Journal of Financial Economics}, 78(1), 2005.

\bibitem{jeannin2008}
M.~Jeannin, G.~Iori, D.~Samuel, ``Modeling Stock Pinning,'' \textit{Quantitative Finance}, 8(8), 2008.

\bibitem{golez2012}
B.~Golez and J.~C.~Jackwerth, ``Pinning in the S\&P~500 Futures,'' \textit{Journal of Financial Economics}, 106(3), 2012.

\bibitem{barbon2021}
A.~Barbon and A.~Buraschi, ``Gamma Fragility,'' Working Paper, University of St.~Gallen, 2021.

\bibitem{filippou2022}
I.~Filippou, P.~A.~Garcia-Ares, F.~Zapatero, ``No Max Pain, No Max Gain: A Case of Predictable Reversal,'' SSRN Working Paper, 2022.

\bibitem{chung2013}
K.~Chung et~al., ``Weekly Options on Stock Pinning,'' Working Paper, William Paterson University, 2013.

\bibitem{ni2021}
S.~X.~Ni, N.~D.~Pearson, A.~M.~Poteshman, J.~S.~White, ``Does Option Trading Have a Pervasive Impact on Underlying Stock Prices?'' \textit{Review of Financial Studies}, 34(4), 2021.

\bibitem{elms2026dw}
N.~Elms, ``Discriminator Ladder Learning: Automatic Semantic Routing for Structured Data Retrieval,'' 2026.

\bibitem{elms2026five}
N.~Elms, ``Five Consequences of One Algorithm,'' 2026.

\end{thebibliography}

\end{document}
